%!TEX ROOT = thesis.tex
\chapter{Introduction}

\section{Overview of Cognitive Impairment, causes, prevalence, and impacts}
Cognitive impairment (\acrshort{CI}) is a medical condition described as a limitation of the human brain to function normally. CI could limit the patient's ability to think, remember, pay attention, communicate, solve problems, and make decisions. The severity of CI can progress through different stages, from mild cognitive impairment (\acrshort{MCI}) to a severe stage known as dementia. MCI is a state in which the patient develops clear signs of cognitive decline but does not require assistance with daily activities \cite{Knopman2014Mild}. Dementia refers to a severe decline in cognitive functions that makes patients rely on others to perform their daily tasks. From a biological point of view, MCI is defined by minor alterations in the functioning of the brain that do not yet result in significant shrinkage or death of brain tissue, which is typical of dementia.

There are various factors and conditions that could lead to CI, including neurodegenerative diseases, vascular disorders, traumatic brain injuries \cite{Nordström2018Traumatic}, genetic factors \cite{Loy2014Genetics}, psychiatric disorders, and other medical conditions. Alzheimer’s disease (\acrshort{AD}) \cite{Scheltens2021Alzheimer} is the most common cause of dementia, contributing to 60\% to 80\% of cases \cite{Alzheimer2024}. AD is caused by the accumulation of amyloid-beta and tau protein and the formation of neurofibrillary tangles. AD leads to degradation in most cognitive abilities, but memory is the most affected function. Vascular dementia (\acrshort{VaD}) is the second most common type of dementia and is caused by a series of small strokes or other conditions that hinder the blood supply to the brain. Some patients may also develop a combination of AD and VaD known as mixed dementia. Other common neurodegenerative diseases are dementia with Lewy body (\acrshort{DLB}) and frontotemporal dementia (\acrshort{FTD}), which also cause CI through different mechanisms. DLB is caused by the accumulation of abnormal protein deposits. Patients with DLB usually experience loss of attention, visual hallucinations, and motor difficulties. FTD is related to degeneration of the frontal and temporal lobes of the brain. The primary symptoms of FTD are changes in personality, behavior, and language abilities.

Furthermore, diseases such as cardiovascular disease, diabetes, and sleep apnea are also related to the development of dementia \cite{Alzheimer2024}. In addition, depression and other psychiatric disorders cause cognitive changes that can mimic early dementia. Obesity, smoking, excessive alcohol consumption, and sedentary lifestyles may also increase the risk of developing dementia-related disorders \cite{WHO2025Dementia}. This highlights the fact that CI is diverse and can have multiple causes.

Dementia-related disorders are among the major health issues in different parts of the world. According to the World Health Organization (\acrshort{WHO}) \cite{WHO2025Dementia}, dementia is one of the leading global causes of death and dependency among older adults. The WHO estimated the prevalence of dementia to be about 55 million cases in 2019. This figure is expected to reach 78 million in 2030 and 139 million in 2050, with 68\% of cases residing in low- and middle-income countries (\acrshort{LMICs}) \cite{Alzheimer2025}. The increasing incidence is attributed to the increase in the aging population because the probability of dementia increases dramatically with age. 

The WHO estimated the global cost of dementia in 2019 at 1.3 trillion US dollars, including the costs of unpaid informal caregivers \cite{WHO2025Dementia}. With the increase in dementia care costs, the total cost is projected to increase to 2.8 trillion US dollars by 2030. This projection calls for immediate changes in healthcare systems and national budgets globally. In addition to economic cost, dementia can cause psychological and social impacts on patients, their caregivers, and family members. Loss of independence can make patients feel distressed. Caregivers and family members provide 5 hours of informal care per day \cite{WHO2025Dementia}. Their demanding caregiving roles could lead to stress, depression, and social isolation. This could also result in a reduction in the productivity of the workforce, as caregivers are often required to reduce their work hours or stop working. In addition, a greater demand is placed on community-level social support systems and long-term care facilities. 

These impacts are particularly severe in LMICs. These countries often face numerous difficulties in terms of care and support for dementia patients due to the lack of healthcare resources, facilities, and access to appropriate services and technologies \cite{Mattap2022Economic}. Moreover, sociocultural factors such as stigma and low awareness of dementia can also delay the development of effective care and support systems in such areas \cite{WHO2025Dementia}. This discrepancy underlines the need for more commitment and studies to identify the challenges in developing countries, where the rates of dementia are increasing dramatically. 

\section{Early Prediction and Classification of Cognitive Impairment}
Pathological changes that cause CI start 10 to 20 years before the onset of symptoms \cite{Beason-Held2013Changes}. These changes could be mainly driven by modifiable risk factors, which are associated with 40\% of dementia cases worldwide \cite{Livingston2020Dementia}. These risk factors include conditions such as hypertension, diabetes, depression, and hearing loss, in addition to lifestyle choices such as physical inactivity and smoking. Management of these risks is possible through early and effective interventions. For example, cognitive decline can be delayed or slowed when hypertension and diabetes are properly controlled with medication and diet at an early age.

The transition from a normal cognitive status to the development of severe CI is gradual. It has been established that about 15\% of patients with MCI are at risk of developing dementia in the next two years \cite{Alzheimer2024}. Additionally, patients with MCI might not experience further decline in cognitive functioning or might even reverse to normal cognition \cite{Alzheimer2024}, especially if they are provided with appropriate treatment and care. Therefore, failure to identify the risk of cognitive impairment early increases the possibility of developing severe and irreversible CI, such as dementia. Moreover, early detection of CI without classifying the potential cause of impairment and the main factors driving the risk might not be sufficient. As discussed earlier, CI is a heterogeneous condition and can develop due to various etiologies. The development of effective intervention and treatment plans requires the classification of distinct patterns and profiles of cognitive changes.

There are many advantages that can be gained from early identification of CI for the individual, caregivers, and society in general. Ideally, identifying the problem can lead to its treatment or at least its rate of progression being reduced so that the person’s cognitive functioning deteriorates at a slower rate and their functional capabilities are preserved longer \cite{Mukadam2019Population}. Therefore, it improves the life experience of the person in need of care and relieves stress on caregivers and health facilities. Furthermore, early detection of CI enables patients and their families to prepare appropriately for future needs in terms of care and finances. It also gives them enough time to prepare for the psychological and social effects caused by loss of cognitive ability. From a larger perspective, early detection can slow down the growing prevalence of dementia, which could save the cost of hospitalizations and long-term care services \cite{Mukadam2019Population}.

Despite these advantages of early detection of CI, the percentage of patients diagnosed early is still low \cite{White2021Observational}. This could be related to various factors such as low awareness, poor or costly availability of screening tools, and the non-specific nature of the early stages of cognitive deterioration. Moreover, the diagnosis of CI is commonly performed based on clinical cognitive assessments of the patient and more detailed neuropsychological tests to establish the diagnosis. This includes interviewing patients, physical examinations, neuropsychological assessments, and possibly imaging or laboratory examinations. This lengthy diagnostic process requires sophisticated equipment and professional medical personnel. In addition, it is usually performed when the patient has already developed irreversible neurological deficits. As a result, current diagnostic methods are costly and ineffective for the early diagnosis of CI.

Widely used screening tools such as the Mini-Mental State Examination (\acrshort{MMSE}) and the Montreal Cognitive Assessment (\acrshort{MoCA}) are not capable of effectively identifying small changes in cognition that are common in the early stages of dementia \cite{Fraser2021Measuring}. Furthermore, MMSE is sensitive to educational levels and cannot be administered to illiterate individuals, as it consists of writing and reading tasks. Functional Magnetic Resonance Imaging (\acrshort{fMRI}) and Positron Emission Tomography (\acrshort{PET}) can help to assess structural and functional changes in the brains of patients with dementia. However, these imaging biomarkers are not ideal in the initial stages of the disease because it is not possible to detect subtle changes in the structure of the brain. Moreover, access to neuroimaging equipment and professional radiologists is costly and restricted, especially in developing countries. Genetics and other biomarkers, including blood biomarkers and cerebrospinal fluid measures, could increase the general accuracy of the diagnostic. However, they involve more detailed procedures and require specialized professional knowledge to interpret. 

The economic and social barriers and the inefficiency of the existing methods prevent a timely diagnosis of CI, particularly in LMICs. Thus, it is necessary to remove these barriers by increasing the level of education, developing accessible screening methods, and improving the healthcare system so that subjects with CI might receive a timely diagnosis and subsequent treatment. During the past few decades, researchers have been exploring the use of machine learning as a key technology to develop more effective solutions for the early prediction of CI risk.

\section{Summary of Cognitive Impairment Diagnosis Using Machine Learning}
Machine learning (\acrshort{ML}) has made outstanding advances in the diagnosis of CI. It offers accuracy and efficiency for early detection \cite{Habehh2021Machine, Li2022Application}. Unlike traditional diagnosis methods, ML models can analyze large amounts of data from different sources. They can find small patterns and early signs of dementia that might be missed. Researchers have explored various ways to develop CI prediction models. They use different types of data from multiple datasets \cite{Javeed2023Machine}. Neuroimaging data has been used extensively to predict CI \cite{Borchert2023Artificial} and to classify dementia subtypes, such as AD, FTD, and VaD \cite{Wang2024Understanding}. To make diagnosis more accurate, researchers in this area have used different neural network designs, steps to prepare the data, and methods to combine different data types \cite{Odusami2024Machine}. They also explored more advanced deep learning (\acrshort{DL}) methods. Examples are Convolutional Neural Networks (CNNs), autoencoders, transformers, and transfer learning \cite{Kaur2024Systematic, Kaur2024Systematic}.

In addition to neuroimaging, other data types related to cognitive assessments are used to predict CI. Digital biomarkers such as speech, language, facial expressions, and handwriting have been explored as automated techniques to assess cognitive functions \cite{Ding2022Digital}. In this area of research, high detection performance is obtained by using natural language processing techniques and DL models \cite{Shi2023Speech}. Neuropsychological assessments and cognitive test scores are also frequently used as non-invasive predictors \cite{Veneziani2024Applications}. Models such as support vector machines (\acrshort{SVM}) and random forest are commonly trained on these data. Combining basic cognitive scores, such as MMSE, with other clinical data has been found to improve prediction accuracy \cite{Basta2023Personalized}. The features of automated tests have also been used to classify the subtypes of dementia. Examples are the digital Clock Drawing Test (\acrshort{dCDT}) \cite{Davoudi2021Classifying} and speech analysis \cite{Yamada2022Speech}. These features helped to differentiate between conditions such as AD, VaD, and DLB.

Machine learning models that use neuroimaging data and cognitive assessments are capable of accurately predicting CI. The accuracy can range from 80\% up to 99\%. However, challenges such as dataset availability and the need for trained professionals and special equipment are limiting their applications. Therefore, affordable and accessible screening tools are crucial in resources-constrained settings to predict CI on a wide scale and reduce dementia prevalence \cite{Magklara2018Current}.

To offer more cost-effective CI diagnosis solutions, researchers have started using readily accessible data such as demographics, lifestyle, medical history, and administrative data. Recent studies have utilized data from longitudinal studies and electronic health records (\acrshort{EHR}) to develop models that can accurately identify early CI. Data from large-scale cohorts, such as the National Alzheimer's Coordinating Center (\acrshort{NACC}), are frequently used in this domain \cite{abbas2024Explainable, Ren2023Improving}. The type of data utilized often includes demographics, medical history, behavioral surveys, and basic functional assessments \cite{Ren2023Improving, Ren2025Identifying}. For these types of structured data, the use of traditional ML algorithms is more common than the use of complex DL models \cite{Zhu2020Machine}. Models such as logistic regression, random forest, and various gradient boosting methods are widely evaluated \cite{Akter2025Using, Cui2024Prediction, Georgiev2024Predicting}.

A significant trend in research is the use of these types of accessible data to predict the future onset of dementia. There are studies predicting risk up to 5, 10, or 13 years in advance \cite{Georgiev2024Predicting, Li2023Early}. Other studies focus on classifying the current cognitive stage. They often use hierarchical binary models. This is to separate normal cases from MCI or MCI from dementia \cite{Ren2023Improving, Yue2023Explainable}. The performance of these models shows great potential for cost-effective CI screening. The classification models achieved area under the curve (\acrshort{AUC}) values ranging from 0.80 to 0.96 \cite{Cui2024Prediction, Richter2025Cognitive}. Even long-term 5-year risk prediction models scored high levels of AUCs between 0.85 and 0.88 \cite{Georgiev2024Predicting, Li2023Early}. Routine clinical data are also combined with other types of data, such as neuroimaging and neuropsychological tests, to classify specific dementia subtypes \cite{Silvia2023Differential, Xue2024AI}. However, the classification of dementia subtypes using only models based on EHR data is still limited to a few studies \cite{Sharma2024Leveraging}. Most of the models that classify the dementia subtypes are based on neuroimaging data \cite{Borchert2023Artificial}.

To support clinical adoption and planning for interventions, the black box nature of these ML models must be addressed. High accuracy alone is not sufficient. For real-world clinical use, trust, transparency, and the ability to generate actionable insights are also required \cite{Martin2023Interpretable}. Therefore, Explainable Artificial Intelligence (\acrshort{XAI}) methods have been integrated into recent CI prediction studies \cite{Martin2023Interpretable}. The most common XAI method found in the literature is the SHapley Additive exPlanations (\acrshort{SHAP}). It is typically applied at the model level. This is done to identify and rank the most important features that influence the prediction of the model, such as age, stroke, or depression \cite{Georgiev2024Predicting, Hossain2025Optimizing, Li2023Early}. Other methods are also used, like Local Interpretable Model-agnostic Explanations (\acrshort{LIME}). These provide patient-level explanations for individual predictions \cite{abbas2024Explainable, Mahamud2025Enhancing, Vyas2022Identifying}. In all of these studies, the XAI methods have consistently identified key predictors. Among the important features are age, cognitive assessment scores (like MMSE), and mental health variables (like depression) \cite{Akter2025Using, Basta2023Personalized, Twait2023Dementia}. Although existing research provides a foundation for cost-effective CI prediction using accessible data, critical challenges remain. This study aims to address these limitations to produce models suitable for wide clinical applications.

\section{Problem Statement}
The early prediction of the risk of CI and the accurate classification of its cause are important. These actions help to control the prevalence of CI and also reduce its cost by supporting the implementation of timely and effective intervention programs \cite{Panegyres2016Early, Rasmussen2019Alzheimers}. Nevertheless, the World Alzheimer’s Report 2021 showed a serious problem. It revealed that 75\% of dementia cases worldwide and 90\% in LMICs are not diagnosed \cite{Gauthier2021World}. According to the report, people with dementia and their caregivers face diagnosis challenges related to lack of access to trained clinicians, fear of diagnosis, and cost of diagnosis. Furthermore, the majority of the cognitive status classification and future CI risk prediction ML models are trained in brain scans and neuropsychological tests \cite{Kaur2024Systematic}. These types of data are expensive and not widely accessible, especially in LMICs and other areas with limited resources. The same case applies to CI subtype classification models, which are mainly developed using neuroimaging data \cite{Wang2024Understanding}.

On the other hand, ML models that use accessible data have critical limitations. Most of these models are developed without validation on external data. This raises serious questions about the generalizability of the models to diverse populations. Moreover, the evaluation of model computation efficiency and the assessment of model calibration and clinical utility are often not reported. Furthermore, training data are typically imbalanced and the class overlap is often not measured and effectively handled in a way that reduces the impacts on the model's performance. Lastly, most existing models do not have the actionable and personalized explanations needed for clinical trust. The explanations are often limited to global feature rankings. They do not provide the counterfactual and patient-specific insights that are needed to guide interventions. All these limitations make existing solutions less reliable and practical \hl{for real clinical use}, particularly in rural areas and LMICs.

\section{Research Question}
How can a methodological machine learning framework overcome the challenges of class overlap and generalizability in developing cost-effective and explainable models for the early prediction of cognitive impairment using accessible data?
This research aims to answer the following key questions:
\begin{itemize}
    \item To what extent can cognitive status be classified at the current visit using cost-effective machine learning models?
    \item How accurately can cognitive impairment risk be predicted five years in advance using cost-effective machine learning models?
    \item How accurately can Alzheimer’s disease be distinguished from other cognitive impairment subtypes?
    \item How can explainable machine learning techniques generate patient-specific insights for cognitive impairment prevention?
\end{itemize}

\section{Research Objectives}
The main objectives of this study are:
\begin{itemize}
    \item To classify current cognitive status (normal cognition, cognitive impairment but not dementia, or dementia) using cost-effective machine learning models.
    \item To predict cognitive status (normal cognition, cognitive impairment but not dementia, or dementia) five years in the future using cost-effective machine learning models.
    \item To classify cognitive impairment subtypes as Alzheimer's disease or non-Alzheimer's disease.
    \item To identify important predictors and generate patient-specific counterfactual explanations for cognitive impairment prevention.
\end{itemize}

\section{Hypothesis}
Based on the research questions, the following hypotheses are proposed:
\begin{itemize}
    \item It is hypothesized that high performance can be achieved in classifying current cognitive status using only accessible data when class overlap is effectively handled. This hypothesis is supported by the theory that cognitive decline progresses gradually, which naturally creates overlapping boundaries between different cognitive stages in the data. Therefore, advanced data preprocessing and modeling techniques must be utilized to resolve these overlapping boundaries.
    \item It is hypothesized that despite the increased complexity and time gap, cost-effective models trained on accessible data from the baseline visit can achieve a clinically acceptable level of accuracy in a five-year future risk prediction task. This assumption is based on the clinical theory that severe cognitive impairment is developed as a progressive condition in which early indicators and symptoms can be observed by analyzing accessible data several years before a formal diagnosis can be established.
    \item It is hypothesized that the classification of cognitive impairment subtypes will be the most challenging task. This is due to the significant similarity of symptoms between Alzheimer’s disease and other subtypes and the low discriminative power of the accessible data. As a result, the use of more predictive features, such as cognitive test scores, might still be needed to increase classification performance.
    \item It is hypothesized that patient-specific model explanations can be developed. The effectiveness of the generated explanation depends on the features on which the model is trained.
\end{itemize}

\section{Key Contributions}
This research makes several key contributions to the body of the research:
\begin{itemize}
    \item It uses demographic, medical history, and behavioral and functional assessment data from thousands of subjects in the NACC study. These data can be easily collected from subjects or their caregivers at minimal or no cost, making the developed model highly accessible.
    \item It proposes ensemble feature selection and hybrid data resampling methods guided by class overlap measures. These systematic methods support the development of cost-effective and generalizable models.
    \item It introduces a novel algorithm for automatic and seamless model selection based on the performance and efficiency of the models. The algorithm promotes the development of cost-effective models.
    \item It performs external validation using data from two different cohorts to assess the generalizability of the models in different healthcare settings.
    \item It assesses the reliability of the models calibration and the models' clinical utility. 
    \item It evaluates an XAI capability to interpret model outputs and counterfactual explanations that can be utilized for personalized risk prevention plans.
\end{itemize}
\vspace{2em}

These contributions are in line with the Sustainable Development Goals (\acrshort{SDGs}) of the United Nations. The research outcome directly supports SDG 3 (Good Health and Well-being). This is achieved by proposing cost-effective models for the early prediction of CI that are accessible and accurate (a formal definition of a cost-effective model is provided in the Glossary, Table~\ref{table:glossary}). In addition, the research work contributes to SDG 10 (Reduced Inequalities) by providing opportunities for equal access to CI prediction in resource-limited settings. Consequently, the healthcare diagnostic workflow could be positively impacted by these contributions. Using accessible data and automated models as an early screening tool, the time and resources required for early diagnosis could be reduced, and unnecessary advanced medical tests might be avoided. As a result, preliminary assessments might be simplified and widely accessible.

The remainder of this thesis is organized into several chapters. Chapter 2 provides a review of the literature, discusses related work, and highlights research gaps. Chapter 3 describes the research methodology, including the proposed methods for data preprocessing, model fine-tuning, external validation, and explanation. Chapter 4 presents and critically analyzes the results of the research work. Chapter 5 discusses the implications of the results with regard to the research objectives. Finally, Chapter 6 concludes the thesis by summarizing the key findings and their significance and highlighting limitations and future work.